\subsection{Phase 4: Validation \& Deployment}

\subsubsection{Step 6: Full Deployment}

The full deployment phase applied all the trained models to annotate the complete CCF Database of 9.2 million sentences across 266,271 articles. We implemented a hierarchical classification strategy that directly mirrors the annotation protocol's structure of primary categories and their associated sub-categories, as illustrated in Figure~\ref{fig:hierarchical_annotation}. 

The hierarchical approach reflects the conceptual organization of our annotation framework with eleven primary detection categories, each with specific sub-categories. The eight thematic frames (\emph{Economic}, \emph{Health}, \emph{Security}, \emph{Justice}, \emph{Political}, \emph{Scientific}, \emph{Environmental}, and \emph{Cultural Frame Detection}) function as primary categories with their internal distinctions—for example, \emph{Economic Frame Detection}, when positive (=1), triggers six economic sub-models (\emph{Negative/Positive Economic Impacts}, \emph{Costs/Benefits of Climate Action}, \emph{Economic Sector Footprint}), while a negative detection (=0) bypasses these entirely. 

\begin{figure}[b!]
\centering
\resizebox{0.95\textwidth}{!}{%
\begin{tikzpicture}[
    % Style definitions with reduced widths
    stage/.style={rectangle, rounded corners=6pt, minimum height=2.3cm, text width=3.8cm, align=center, font=\small},
    input/.style={stage, draw=orange!70, fill=orange!5, text width=3.5cm},
    detect/.style={stage, draw=blue!70, fill=blue!5, text width=3.8cm},
    process/.style={stage, draw=green!70, fill=green!5, text width=3.5cm},
    skip/.style={stage, draw=red!70, fill=red!5, text width=3.5cm},
    storage/.style={cylinder, draw=purple!70, fill=purple!5, shape border rotate=90, minimum height=2.8cm, minimum width=3cm, align=center, font=\small, aspect=0.3},
    arrow/.style={->, >=latex, line width=1.5pt, color=gray!60},
    yes/.style={arrow, color=green!60},
    no/.style={arrow, color=red!60},
    label/.style={font=\tiny\bfseries}
]

% Stage 1: Input
\node[input] (input) {
    \scriptsize\textbf{Input Layer}\\[-2pt]
    \scriptsize\textbf{Sentence Processing}\\[-2pt]
    \textit{\scriptsize "Carbon pricing will}\\[-2.5pt]
    \textit{\scriptsize boost clean energy jobs"}\\[-2pt]
    \scriptsize{One of 9.2M sentences}
};

% Stage 2: Detection (parallel branches)
\node[detect, above right=1.2cm and 1.2cm of input] (detect_yes) {
    \scriptsize\textbf{Detection Layer}\\[-2pt]
    \scriptsize\textbf{Main Category}\\[-2pt]
    \textit{\scriptsize e.g., Economic Frame}\\[-2pt]
    \colorbox{green!20}{\scriptsize\textbf{Prediction = 1}}\\[-2pt]
    \scriptsize{→ Activate sub-models}
};

\node[detect, below right=1.2cm and 1.2cm of input] (detect_no) {
    \scriptsize\textbf{Detection Layer}\\[-2pt]
    \scriptsize\textbf{Main Category}\\[-2pt]
    \textit{\scriptsize e.g., Health Frame}\\[-2pt]
    \colorbox{red!20}{\scriptsize\textbf{Prediction = 0}}\\[-2pt]
    \scriptsize{→ Skip sub-models}
};

% Stage 3: Conditional execution
\node[process, right=1.5cm of detect_yes] (execute) {
    \scriptsize\textbf{Sub-classification}\\[-2pt]
    \scriptsize\textbf{Apply Sub-models}\\[-2pt]
    \textit{\scriptsize Benefits: 1, Costs: 0}\\[-2.5pt]
    \textit{\scriptsize Positive: 1, Negative: 0}\\[-2pt]
    \scriptsize{Sub-models executed}
};

\node[skip, right=1.5cm of detect_no] (bypass) {
    \scriptsize\textbf{Skip Processing}\\[-2pt]
    \scriptsize\textbf{No Sub-models}\\[-2pt]
    \textit{\scriptsize Sub-categories}\\[-2.5pt]
    \textit{\scriptsize not needed}
};

% Stage 4: Storage - positioned at center between execute and bypass
\node[storage] (database) at ($(execute.east)!0.5!(bypass.east) + (2cm,0)$) {
    \scriptsize\textbf{CCF Database}\\[-2pt]
    \scriptsize 9.2M sentences\\[-2.5pt]
    \scriptsize 266K articles
};

% Connections
\draw[arrow] (input.15) -- (detect_yes.west);
\draw[arrow] (input.-15) -- (detect_no.west);

\draw[yes, line width=1.8pt] (detect_yes) -- node[above, label] {\textcolor{green!70}{if = 1}} (execute);
\draw[no, line width=1.8pt] (detect_no) -- node[below, label] {\textcolor{red!70}{if = 0}} (bypass);

\draw[arrow] (execute.east) -- (database.west);
\draw[arrow] (bypass.east) -- (database.west);

% Architecture info box
\node[below=0.3cm of input, draw=gray!50, fill=gray!3, rounded corners=4pt, text width=3.2cm, align=center, font=\tiny] {
    \textbf{Parallel Processing}\\
    65 detection models\\
    BERT / CamemBERT
};

% Visual indicator for parallel processing - better positioned
\node[above=0.3cm of detect_no, font=\tiny\itshape, text=gray!50] {61 other models};
\node[above=0.1cm of detect_no, font=\tiny, text=gray!50] {$\vdots$};

\end{tikzpicture}
}
\caption{Hierarchical annotation pipeline. The system processes each sentence through detection models that determine whether corresponding sub-category models should be applied, enabling efficient annotation of the entire corpus.}
\label{fig:hierarchical_annotation}
\end{figure}

Similarly, the three other primary categories follow this conditional logic: \emph{Actors/Messengers Detection} activates nine actor-type sub-models (\emph{Health Expert}, \emph{Economic Expert}, \emph{Security Expert}, \emph{Legal Expert}, \emph{Cultural Expert}, \emph{Natural Scientist}, \emph{Social Scientist}, \emph{Activist}, \emph{Public Official}), \emph{Event Detection} triggers eight event-type classifications (\emph{Extreme Weather Event}, \emph{Meeting/Conference}, \emph{Publication}, \emph{Election}, \emph{Policy Announcement}, \emph{Judiciary Decision}, \emph{Cultural Event}, \emph{Protest}), and \emph{Solutions Detection} applies two solution sub-categories (\emph{Mitigation Strategy}, \emph{Adaptation Strategy}). The deployment also processed standalone primary categories without hierarchical structure—\emph{Emotional Tone} (\emph{Positive Emotion}/\emph{Negative Emotion}/\emph{Neutral Emotion}), \emph{Geographic Focus} (\emph{Canadian Context}), \emph{Urgency/Alarmism}. 

In addition to the custom-trained classification models, we employed pre-trained Named Entity Recognition (NER) models, selected based on empirical evaluation for optimal performance on our dataset. We implemented a hybrid approach tailored to each language: for English texts, we used BERT-base-NER\footnote{Available at: \url{https://huggingface.co/dslim/bert-base-NER}} for all three entity types (Person, Organization, Location), while for French texts, we combined spaCy's fr\_core\_news\_lg model\footnote{Available at: \url{https://spacy.io/models/fr}} for person entities with CamemBERT-NER\footnote{Available at: \url{https://huggingface.co/Jean-Baptiste/camembert-ner}} for organization and location entities. This hybrid strategy was used to leverage the respective strengths of each model: spaCy's superior performance on French person names and CamemBERT-NER's robust identification of French organizational and geographical entities (see Table~\ref{tab:ner_performance} in the Appendix for performance metrics).

\subsubsection{Step 7 \& 8: Final validation and inter-coder reliability}

The final validation phase employs a stratified sampling scheme designed to ensure robust performance evaluation across all 65 annotation categories. Following the completion of full database annotation (9.2 million sentences), we extracted a balanced validation sample of 1,000 sentences (500 per language) using root-inverse probability weighting\footnote{The weighting formula $w_i \propto \sum_j 1/\sqrt{f_j}$ assigns higher sampling probability to sentences containing rare categories, where $f_j$ represents the frequency of category $j$ present in sentence $i$. For example, a sentence annotated with both \textit{Economic Frame} (appearing in 15.4\% of the corpus) and \textit{Loss of Indigenous Practices} (appearing in 0.28\% of the corpus) would receive a weight of approximately $1/\sqrt{0.154} + 1/\sqrt{0.0028} = 2.55 + 18.90 = 21.45$, making it roughly 8 times more likely to be selected than a sentence containing only \textit{Economic Frame}.} combined with strict constraints to ensure statistical validity. Two critical constraints were used : first, a minimum threshold of 40 positive examples per category to ensure sufficient statistical power for reliable performance estimation. Without this constraint, rare categories like \textit{Military Base Disruption} (0.00004\% prevalence in the database) would have too few examples to meaningfully assess model accuracy. Second, a 35\% maximum allocation to prevent any single prevalent category from monopolizing the validation set. For instance, without this constraint, \textit{Actors/Messengers Detection} (present in 48.8\% of database sentences) could theoretically dominate the entire sample, leaving insufficient space for other important categories.

The validation results demonstrate robust performance across all annotation dimensions, with an overall F1 macro score of 0.866 (0.869 for English, 0.864 for French). As detailed in Table~\ref{tab:detailed_validation_metrics} (Appendix B), primary detection categories achieve the highest performance, with \textit{Canadian Context} (F1=0.982) and \textit{Actors/Messengers Detection} (F1=0.961) showing near-perfect classification. Thematic frames maintain strong performance despite their semantic complexity, with \textit{Scientific Frame Detection} (F1=0.890) and \textit{Environmental Frame Detection} (F1=0.871) leading this category. The hierarchical classification strategy proves particularly effective, as evidenced by consistently higher recall scores for detection categories (mean recall=0.891) that trigger conditional sub-category evaluation. Notably, rare categories such as \textit{Military Base Disruption} and \textit{Military Disaster Response} show perfect precision despite minimal representation.

To assess temporal stability, we conducted stratified validation across five temporal periods (1980s-2020s) that revealed minimal performance drift ($\Delta$F1 < 1.1\%) despite evolving media discourse patterns. This temporal consistency validates the models' generalization capacity across the full 46-year corpus span. Inter-coder reliability assessment, currently in progress with a second independent annotator, will evaluate the same 1000 randomly selected sentences from the validation set. Final inter-coder metrics will be reported upon completion of the dual-annotation process.

\begin{table}[h!]
\centering
\caption{Overall validation performance metrics}
\label{tab:final_validation_metrics}
\begin{tabular}{lccc}
\toprule
\textbf{Language} & \textbf{F1 Macro} & \textbf{F1 Micro} & \textbf{F1 Weighted} \\
\midrule
English (EN) & 0.869 & 0.909 & 0.911 \\
French (FR) & 0.864 & 0.902 & 0.905 \\
\midrule
\textbf{Combined (ALL)} & \textbf{0.866} & \textbf{0.905} & \textbf{0.908} \\
\bottomrule
\end{tabular}
\end{table}

\section{Appendix B: Complete Framework and Performance Tables}
\label{sec:appendix_tables}
\setcounter{table}{0}  % Reset table counter for appendix
\renewcommand{\thetable}{B\arabic{table}}  % Prefix appendix tables with 'B'
\renewcommand\theHtable{AppendixB.\thetable}  % Fix hyperref links for appendix tables

\landscape
{\footnotesize
\phantomsection
\begin{longtable}{p{0.5cm}p{5cm}p{16cm}}
\caption{Complete CCF annotation framework: All 65 categories with operational definitions}
\label{tab:complete_framework_appendix}
\\
\toprule
\textbf{\#} & \textbf{Category} & \textbf{Description} \\
\midrule
\endfirsthead
\multicolumn{3}{c}{\tablename\ \thetable\ -- \textit{Continued from previous page}} \\
\toprule
\textbf{\#} & \textbf{Category} & \textbf{Description} \\
\midrule
\endhead
\midrule
\multicolumn{3}{r}{\textit{Continued on next page}} \\
\endfoot
\bottomrule
\endlastfoot

\multicolumn{3}{l}{\cellcolor{gray!10}\textbf{THEMATIC FRAMES}} \\
\midrule
\multicolumn{3}{l}{\textit{Economic Frame}} \\
1 & Economic Frame Detection & Presence of any economic dimension of climate change \\
2 & Negative Economic Impacts & Economic losses from climate change (e.g., floods, infrastructure damage) \\
3 & Positive Economic Impacts & Economic benefits from climate action (e.g., green jobs, energy savings) \\
4 & Costs of Climate Action & Economic costs of transitioning to low-carbon economy \\
5 & Benefits of Climate Action & Economic opportunities from climate investments \\
6 & Economic Sector Footprint & Carbon intensity of economic sectors \\
\midrule
\multicolumn{3}{l}{\textit{Health Frame}} \\
7 & Health Frame Detection & Presence of any relationship between climate and health \\
8 & Negative Health Impacts & Heat stress, disease spread, respiratory issues, mental health burdens, mortality \\
9 & Positive Health Impacts* & Potential benefits to health from climate change \\
10 & Health Co-benefits & Better air quality, improved diets, avoided premature deaths, mental well-being \\
11 & Health Sector Footprint* & Carbon footprint of the healthcare sector \\
\midrule
\multicolumn{3}{l}{\textit{Security Frame}} \\
12 & Security Frame Detection & Presence of any security dimension \\
13 & Military Disaster Response* & Army called in for fires, floods, evacuations or relief \\
14 & Military Base Disruption* & Climate threats to military facilities or readiness \\
15 & Climate-Driven Displacement & Military management of evacuations or refugee camps \\
16 & Resource Conflict & Tensions or violence over water, land or minerals worsened by climate change \\
17 & Defense Sector Footprint* & Carbon footprint of military operations \\
\midrule
\multicolumn{3}{l}{\textit{Justice Frame}} \\
18 & Justice Frame Detection & Presence of any social justice angle \\
19 & Winners \& Losers & Groups that benefit or suffer from climate measures (workers, vulnerable populations) \\
20 & North-South Responsibility & Common-but-differentiated responsibilities between high and low-income countries \\
21 & Unequal Impacts & Disproportionate climate impacts on vulnerable groups \\
22 & Unequal Access & Differential access to resources, adaptation, or clean technology \\
23 & Intergenerational Justice & Responsibilities to future generations \\
\midrule
\multicolumn{3}{l}{\textit{Political Frame}} \\
24 & Political Frame Detection & Presence of any policy measure or political discussion \\
25 & Policy Measures & Concrete climate laws, regulations, or programmes under debate or in force \\
26 & Political Debate & Parliamentary disputes, party platforms \\
27 & Political Positioning & Strategic positioning or partisan signalling around climate \\
28 & Public Opinion & Polls or surveys on climate and energy \\
\midrule
\multicolumn{3}{l}{\textit{Scientific Frame}} \\
29 & Scientific Frame Detection & Presence of any scientific aspect \\
30 & Scientific Controversy & Debates on climate change reality, causes, thresholds, geoengineering ethics \\
31 & Discovery \& Innovation & New findings on climate impacts or emerging technologies (e.g., carbon capture) \\
32 & Scientific Uncertainty & Expressions of doubt or uncertainty about climate science \\
33 & Scientific Certainty & Strong consensus statements about climate science \\
\midrule
\multicolumn{3}{l}{\textit{Environmental Frame}} \\
34 & Environmental Frame Detection & Presence of any ecological dimension \\
35 & Habitat Loss & Destruction of natural habitats due to climate change \\
36 & Species Loss & Extinction or endangerment of species \\
\midrule
\multicolumn{3}{l}{\textit{Cultural Frame}} \\
37 & Cultural Frame Detection & Presence of any cultural aspect \\
38 & Artistic Representation & Books, documentaries, plays, exhibitions portraying climate themes \\
39 & Event Disruption & Sports or cultural events threatened or cancelled due to climate conditions \\
40 & Loss of Indigenous Practices & Erosion of traditional hunting, fishing, or cultural rituals linked to climate \\
41 & Cultural Sector Footprint & Emissions from film production, fashion, large festivals, etc. \\
\midrule
\multicolumn{3}{l}{\cellcolor{gray!10}\textbf{PRIMARY CATEGORIES}} \\
\midrule
\multicolumn{3}{l}{\textit{Actors/Messengers}} \\
42 & Actors/Messengers Detection & Presence of any messenger, expert or authority figure \\
43 & Health Expert & Physicians, epidemiologists, health ministers, public health officials \\
44 & Economic Expert & Economists, finance ministers, market analysts, central bank officials \\
45 & Security Expert & Military officers, defense strategists, security scholars \\
46 & Legal Expert & Lawyers, judges, legal scholars, justice ministers \\
47 & Cultural Expert & Artists, writers, athletes, arts scholars commenting on climate change \\
48 & Natural Scientist & Natural science researchers or academics \\
49 & Social Scientist & Social science researchers or academics \\
50 & Activist & Environmental NGO spokespeople or well-known climate activists \\
51 & Public Official & Civil servants or government officials speaking in an official capacity \\
\midrule
\multicolumn{3}{l}{\textit{Events}} \\
52 & Event Detection & Mentions at least one specific event type \\
53 & Extreme Weather Event & Arrival or unfolding of floods, wildfires, hurricanes, heatwaves, etc. \\
54 & Meeting/Conference & International meetings such as COP, UN summits, major national conferences \\
55 & Publication & Publication of governmental, NGO or scientific reports (e.g., IPCC, Lancet) \\
56 & Election & Climate issues raised during local, provincial or national elections \\
57 & Policy Announcement & Debut or unveiling of new climate laws, regulations or action plans \\
58 & Judiciary Decision & Court rulings, trials, regulatory hearings on climate or environment \\
59 & Cultural Event & Sports, artistic, or cultural events (Olympics, marathon, concert, theatre) \\
60 & Protest & Demonstrations or strikes (e.g., climate strike, anti-pipeline protest) \\
\midrule
\multicolumn{3}{l}{\textit{Solutions}} \\
61 & Solutions Detection & Mentions any mitigation or adaptation measure \\
62 & Mitigation Strategy & Measures to reduce GHG emissions or enhance carbon sinks \\
63 & Adaptation Strategy & Measures to increase social or ecological resilience to climate impacts \\
\midrule
\multicolumn{3}{l}{\cellcolor{gray!10}\textbf{EMOTIONAL TONE}} \\
64 & Positive Emotion & Hope, optimism, enthusiasm about climate solutions \\
65 & Negative Emotion & Fear, anxiety, despair about climate impacts \\
66 & Neutral Emotion & No clear emotional tone, factual reporting \\
\midrule
\multicolumn{3}{l}{\cellcolor{gray!10}\textbf{GEOGRAPHIC FOCUS}} \\
67 & Canadian Context & Canadian places, actors, data, and policies \\
\midrule
\multicolumn{3}{l}{\cellcolor{gray!10}\textbf{URGENCY/ALARMISM}} \\
68 & Urgency/Alarmism & Conveys immediate danger, crisis, or ``code red'' urgency \\
\bottomrule
\end{longtable}
\vspace{0.5em}
\noindent\footnotesize
* Insufficient training data (fewer than 10 positive examples in training set).
} % end \footnotesize
\endlandscape